\subsection{分次模的分式模}

设 A 是分次环，M 是分次 A-模，\(\Delta\) 是次数幺半群；在本节中我们假设 \(\Delta\) 是一个群。回忆（《代数学》第 II 章，§ 11），A 与 M 分别是加群的直和

\[
\mathbf {A} = \bigoplus_ {i \in \Delta} \mathbf {A} _ {i}, \quad \mathbf {M} = \bigoplus_ {i \in \Delta} \mathbf {M} _ {i}
\]

其中 \(\mathbf{A}_i\mathbf{A}_j\subset \mathbf{A}_{i + j}\) 与 \(\mathrm{A}_i\mathrm{M}_j\subset \mathrm{M}_{i + j}\) 对所有 \(i,j\in \Delta\) 成立。设 S 是 A 的一个所有元素都是齐次的乘性子集。对所有 \(i\in \Delta\)，我们记 \(\mathbf{S}_i = \mathbf{S}\cap \mathbf{A}_i\)，并用 \((\mathrm{S}^{-1}\mathbf{M})_i\) 表示元素 \(m^{\prime}\)（\(\mathbf{S}^{-1}\mathbf{M}\) 的）的集合，对它们存在元素 \(j\)、\(k\)（属于 \(\Delta\)），元素 \(m\in \mathbf{M}_j\) 与元素 \(s\in \mathbf{S}_k\)，使得 \(j - k = i\) 且 \(m^{\prime} = m / s\)。 若 \((m_q')_{1\leqslant q\leqslant r}\) 是 \(\mathbf{S}^{-1}\mathbf{M}\) 中元素的有限族，使得

\[
m _ {q} ^ {\prime} \in (\mathrm{S} ^ {- 1} \mathrm{M}) _ {i (q)},
\]

则存在元素 \(j(q) \in \Delta\) 与 \(k \in \Delta\)、元素 \(m_q \in \mathrm{M}_{j(q)}\)（\(1 \leqslant q \leqslant r\)）以及元素 \(s \in \mathrm{S}_k\)，使得 \(m_q' = m_q / s\) 对 \(1 \leqslant q \leqslant r\) 成立（第 1 节，注 2）。

命题 21. 环 \(\mathbf{S}^{-1}\mathbf{A}\) 连同族 \(((\mathbf{S}^{-1}\mathbf{A})_i)\) 是分次环，而 \(\mathbf{S}^{-1}\mathbf{M}\) 连同族 \(((\mathbf{S}^{-1}\mathbf{M})_i)\) 是分次环 \(\mathbf{S}^{-1}\mathbf{A}\) 上的分次模。典范映射 \(i_{\mathbf{A}}^{\mathbf{S}}\) 与 \(i_{\mathbf{M}}^{\mathbf{S}}\) 是 0 次齐次的。

设 \(m \in M_j\)、\(s \in S_k\)、\(m' \in M_{j'}\)、\(s' \in S_{k'}\)，并设 \(j - k = j' - k' = i\)；则 \((m/s) - (m'/s') = (s'm - sm')/ss'\)，且 \(s'm - sm' \in M_{j+k'} = M_{j'+k}\)，且 \(ss' \in S_{k+k'}\)，因而按定义 \((m/s) - (m'/s') \in (\mathrm{S}^{-1}\mathrm{M})_i\)；这表明诸 \((\mathrm{S}^{-1}\mathrm{M})_i\) 是 \(S^{-1}\mathrm{M}\) 的加法子群。 这些群的和是整个 \(S^{-1}\mathrm{M}\)：因为每个 \(x \in S^{-1}\mathrm{M}\) 都可写成 \(m/s\)，其中 \(m \in \mathbf{M}\)、\(s \in \mathbf{S}\)；由假设 \(s\) 是齐次的，而 \(m\) 是齐次元素 \(m_j\) 之和；于是 \(x\) 是诸 \(m_j / s\) 之和，其中每个都属于某个子群 \((\mathbf{S}^{-1}\mathbf{M})_i\)。 最后，诸 \((\mathbf{S}^{-1}\mathbf{M})_i\) 的和是直和；因为考虑元素的有限族 \(x_q\)（\(1 \leqslant q \leqslant n\)），它满足 \(x_q \in (\mathbf{S}^{-1}\mathbf{M})_{i(q)}\)，其中指标 \(i(q)\) 互不相同，并设 \(\sum_{q=1}^{n} x_q = 0\)。 每个 \(x_q\) 都可写成 \(x_q = m_q / s\)，其中 \(s \in S_k\) 且 \(m_q \in M_{i(q)+k}\)；该假设蕴含存在 \(s' \in S\) 使得 \(s'\left(\sum_{q=1}^{n} m_q\right) = 0\)； 若诸 \(s'm_q\) 不全为零，我们将得到矛盾，因为若 \(s' \in S_d\)，则 \(s'm_q \in M_{i(q)+d}\)，而诸 \(i(q) + d\) 互不相同。由此得出 \(x_q = 0\)，对每个指标 \(q\)。

立即可验证，若 \(a \in (\mathrm{S}^{-1}\mathrm{A})_i\) 且 \(x \in (\mathrm{S}^{-1}\mathrm{M})_j\)，则 \(ax \in (\mathrm{S}^{-1}\mathrm{M})_{i+j}\)。把这个结果应用于 \(\mathbf{M} = \mathbf{A}\) 的情形，我们首先看到 \(\mathrm{S}^{-1}\mathbf{A}\) 是以诸 \((\mathrm{S}^{-1}\mathrm{A})_i\) 分次的环；然后看到 \(\mathrm{S}^{-1}\mathrm{M}\) 是 \(\mathrm{S}^{-1}\mathrm{A}\) 上的分次模。最后，由于 \(1 \in \mathrm{A}_0\)，\(i_{\mathrm{A}}^{\mathrm{S}}\) 与 \(i_{\mathrm{M}}^{\mathrm{S}}\) 是 0 次齐次的。

命题 22. 设 A（相应地，B）是型为 \(\Delta\) 的分次环，M（相应地，N）是分次环 A（相应地，B）上的分次模，S（相应地，T）是 A（相应地，B）的乘性子集且其所有元素都是齐次的，\(f: \mathrm{A} \to \mathrm{B}\) 是 0 次齐次同态且满足 \(f(\mathrm{S}) \subset \mathrm{T}\)，而 \(u: \mathrm{M} \to \mathrm{N}\) 是 \(k\) 次齐次的 A-线性映射。则同态 \(f': \mathrm{S}^{-1}\mathrm{A} \to \mathrm{T}^{-1}\mathrm{B}\)（由 \(f\) 导出；第 1 节，命题 2）是 0 次齐次的，且 \((\mathrm{S}^{-1}\mathrm{A})\)-线性映射

\[
u ^ {\prime} \colon \mathrm{S} ^ {- 1} \mathrm{M} \rightarrow \mathrm{T} ^ {- 1} \mathrm{N}
\]

由 f 与 u 导出的这一映射（第 2 节，命题 5）是 k 次齐次的。

这由等式 \(f'(a / s) = f(a) / f(s)\) 与 \(u'(m / s) = u(m) / f(s)\) 立即得出。

最后我们注意到，若 E 是分次 A-代数而 S 是由齐次元素组成的 A 的乘性子集，则 \(S^{-1}E\) 连同其 \((S^{-1}A)\)-代数结构（第 8 节）与分次 \((S^{-1}E)_{t}\) 是一个分次 \(S^{-1}A\)-代数，这由定义立即得出。

